% ============================================================
%  part2/ch11-relaxation.tex  —  Chapter 11: Why Systems Relax
% ============================================================

\chapter{Why Systems Relax}
\label{ch:relaxation}

\epigraph{%
  Everything flows.%
}{Heraclitus}

\section{Relaxation as free-energy descent}

A recurring theme across physics, chemistry, and biology is that
systems evolve toward lower-energy configurations.
This chapter makes that precise.

\begin{definition}[Stable equilibrium]
  A field $\phi^*$ is a \emph{stable equilibrium} of the
  gradient flow if $\frac{\delta\mathcal{F}}{\delta\phi}\big|_{\phi^*} = 0$
  and the Hessian of $\mathcal{F}$ at $\phi^*$ is positive definite.
\end{definition}

\begin{theorem}[Convergence to equilibrium]
  \label{thm:convergence}
  If $\mathcal{F}$ is bounded below and coercive, gradient flow
  solutions converge to a critical point of $\mathcal{F}$ as
  $t \to \infty$.
\end{theorem}

\begin{StatusBox}{Proven Framework}
  Theorem~\ref{thm:convergence} holds under standard conditions
  on $\mathcal{F}$ (e.g., Łojasiewicz inequality).
  For the coupled RSVP system, global convergence remains
  \statusconjectured{}.
\end{StatusBox}

\section{Metastability and slow relaxation}

Not all systems reach global equilibrium quickly.
Energy barriers can trap systems in metastable states for
long times.

\begin{definition}[Metastable state]
  A field $\phi$ is \emph{metastable} if it is a local minimum
  of $\mathcal{F}$ but not a global minimum.
\end{definition}

\begin{TechNote}[Coarsening as relaxation through metastability]
  The coarsening process studied in Part~III is a sequence of
  escapes from metastable states.
  Each domain configuration is metastable.
  The system slowly evolves toward larger domains, each of
  which is also metastable, until only one domain remains.
  The cosmic web, in the RSVP interpretation, may represent
  a very slowly relaxing metastable configuration.
\end{TechNote}

\begin{HolonomyNote}
  Traversing Part~II as a whole reveals the following invariant:
  every chapter has been about the same dynamics expressed at
  different levels of abstraction.

  Fields, flows, entropy, and variational principles are not
  four different subjects.
  They are four coordinate charts on the same object:
  the gradient descent of a free-energy functional.

  The holonomy of Part~II is the recognition that
  $\frac{d\mathcal{F}}{dt} \leq 0$ is not a theorem about
  a specific system.
  It is a structural consequence of the gradient-flow
  construction that will appear in every subsequent part.
\end{HolonomyNote}

\WhatSurvived{Relaxation $\phi(0) \mapsto \phi(\infty)$}{%
  Equilibrium structure\\
  Conserved quantities\\
  Topological invariants%
}{%
  Initial condition details\\
  Transient structures\\
  Metastable history%
}{%
  Full equilibrium state\\
  recoverable if equations\\
  and conserved quantities\\
  are known.%
}{%
  $\ker$: all initial conditions\\
  in the same basin of\\
  attraction.%
}
